% C.3bp. Final after complete root and two independent peer audits.
\paragraph{C.3bp. Mixed exceptional intersections in every resolved dimension
(Proved reduction).}
We first give an intrinsic intersection assertion over any
algebraically closed field $k$.
Let $Z$ be regular of pure dimension $d+1$, where $d\ge2$, and let
$E=\bigcup_{i=1}^cE_i$ be a projective SNC divisor whose components
are smooth projective $k$-varieties of dimension $d$.
Suppose that
\[
 D=\sum_i a_iE_i,\qquad a_i\in\mathbf Z_{>0},\qquad
                         \mathcal O_E(-D)\ \text{is ample}.
\]
Fix an ample line bundle $H$ on the whole of $E$, and write
$H_i=H|_{E_i}$,
$A_i=\operatorname{Pic}^0_{E_i/k,\mathrm{red}}$,
$B_i=\operatorname{Alb}(E_i)=A_i^\vee$.
For $L_i\in A_i(k)$, put
\[
 z_i(L_i)=c_1(L_i)c_1(H_i)^{d-2}\cap[E_i]
                              \in\mathrm{CH}_1(E_i).
\]
Intersect the pushforward of $\sum_i z_i(L_i)$ into $Z$ with each
$E_j$, and take its degree-zero Albanese class on that component.
This defines a homomorphism
\begin{equation}
       \Lambda_H:\prod_iA_i\longrightarrow\prod_iB_i.
                                                        \tag{MA.1}
\end{equation}
It is an isogeny. This assertion uses neither characteristic-zero
Hodge theory nor a standard conjecture.

Now impose precisely the resolved positive-characteristic
hypotheses of C.3bn: $k$ is algebraically closed of characteristic
$p>0$; $R$ is an excellent henselian local coefficient-$k$ algebra,
with residue field $k$, obtained by strict localization of a
finite-type $k$-scheme; $Z\to\operatorname{Spec}R$ is projective
and regular of pure dimension $d+1$, and is an isomorphism off
the closed point; $E$ is its full reduced closed fibre.
In addition to C.3bn we assume here that $E$ is SNC.
The displayed positive divisor $D$ has full support $E$;
ampleness on $E$ is equivalent to ampleness of
$\mathcal O_D(-D)$ on its nilpotent thickening
\citep[Lemma 30.17.5]{rootStacks2026NilpotentAmple}.
Let $X=Z\setminus E$. Use the C.3bn notation
\[
 S_d(E)=F_\gamma^dK_0(E)_{\mathbf Q},\qquad J_K^d(E),\qquad
 \theta_d:\mathrm{CH}_1(E)_{\mathbf Q}\longrightarrow S_d(E).
\]
Thus $J_K^d(E)$ is generated by endpoint differences of actual
classes in $F_\gamma^dK_0(E\times P)_{\mathbf Q}$ for smooth
connected projective parameter curves $P$.

Define
\begin{align}
 \deg_E:S_d(E)&\longrightarrow\mathbf Q^c,&
 s&\longmapsto
 \bigl(\deg\operatorname{ch}_d(s|_{E_i})\bigr)_i,\notag\\
 S_{d,0}(E)&=\ker(\deg_E),\notag\\
 \alpha_E:S_{d,0}(E)&\longrightarrow
                   \bigoplus_iB_i(k)\otimes\mathbf Q,&
 s&\longmapsto
 \bigl(\operatorname{alb}_{E_i}
             (\operatorname{ch}_d(s|_{E_i}))\bigr)_i,\notag\\
 S_{d,00}(E)&=\ker(\alpha_E).                           \tag{MA.2}
\end{align}
Then the exact remaining quotient is
\begin{equation}
 B^d(X)_{\mathbf Q}\simeq
 \frac{S_{d,00}(E)}
 {(J_K^d(E)\cap S_{d,00}(E))+
  (\theta_d\mathrm{CH}_1(E)_{\mathbf Q}\cap S_{d,00}(E))}.
                                                        \tag{MA.3}
\end{equation}
Moreover, if $U_H$ is the subgroup generated over $\mathbf Q$ by
the classes $\theta_d(z_i(L_i))$, then
\begin{equation}
 U_H\subset J_K^d(E)\cap S_{d,0}(E),\qquad
 \alpha_E(U_H)=\bigoplus_iB_i(k)\otimes\mathbf Q.         \tag{MA.4}
\end{equation}
No existence of the displayed resolution in arbitrary dimension
is asserted, and (MA.3) is not a vanishing statement.

\begin{proof}
\emph{1. Mixed polarizations without a Lefschetz conjecture.}
Let $V$ be smooth connected projective of dimension $n\ge1$ and
let $\gamma\in\mathrm{CH}_1(V)$ be a one-cycle class.
Intersection and the Albanese map give an algebraic homomorphism
\[
 \kappa_\gamma:\operatorname{Pic}^0_{V/k,\mathrm{red}}
           \longrightarrow\operatorname{Alb}(V),\qquad
 L\longmapsto\operatorname{alb}_V(c_1(L)\cap\gamma).
\]
For clarity, its algebraicity requires no assertion about an
arbitrary Chow functor: represent $\gamma$ by a finite signed
sum of integral curves, normalize each curve, and use pullback
on Picard varieties followed by its Jacobian-to-Albanese
pushforward. Proper pushforward and divisor intersection give
the displayed map on points. Rational equivalence of
$\gamma$ does not change that map, and maps of abelian varieties
over an algebraically closed field are determined by their points.
Put $\mu_\gamma=-\kappa_\gamma$.
For a smooth curve $f:C\to V$, let
$r=f^*:\operatorname{Pic}^0_{V/k,\mathrm{red}}\to J_C$.
The sign audited in C.3bo gives
\begin{equation}
                 \mu_{f_*[C]}=r^\vee\lambda_Cr.
                                                        \tag{MA.5}
\end{equation}
Here $\lambda_C$ is the principal polarization of $J_C$.
It follows from
\citet[Lemma 6.9 and Remark 6.10(c), author p.~23]
{D_MilneJacobians2021}; the minus sign is the actual curve
autoduality sign, not a convention changed with the dimension.

If $L_1,\ldots,L_{n-1}$ are ample, write
\[
 \mu_{L_1,\ldots,L_{n-1}}
   =\mu_{\,c_1(L_1)\cdots c_1(L_{n-1})\cap[V]}.
\]
A positive integer multiple of this homomorphism is a polarization.
Here is a complete finite-restriction proof.
Starting with $Y_0=V$, choose successively smooth connected
divisors
\[
 Y_j\in\big|L_j^{\,b_j}|_{Y_{j-1}}\big|,
 \qquad b_j>0,\qquad j=1,\ldots,n-1.
\]
The integers may be chosen in this order and sufficiently large.
At each step $\dim Y_{j-1}\ge2$, and
\begin{equation}
 H^1\bigl(Y_{j-1},L_j^{-b_j}|_{Y_{j-1}}\bigr)
 \simeq
 H^{\dim Y_{j-1}-1}
 \bigl(Y_{j-1},\omega_{Y_{j-1}}\otimes
                        L_j^{b_j}|_{Y_{j-1}}\bigr)^\vee
 =0.                                                    \tag{MA.6}
\end{equation}
We use the smooth proper Serre duality of
\citet[Lemma 48.27.1(5),(7), Lemma 48.27.5(4), and
Remark 48.27.6]{D_Stacks0FVU}, followed by Serre vanishing
\citep[Lemma 30.17.1]{frC3azStacksSerre}.
Smoothness and connectedness of the high sections are supplied
by the Bertini argument in the
\citet[proof of Lemma 45.8.3]{frC3azStacksBertini}.

The divisor sequence and (MA.6) make the restriction on
$H^1(\mathcal O)$ injective at each step. Hence each restriction
of reduced Picard varieties has finite kernel, even if the
unreduced Picard schemes are not smooth. The composite
\[
 r:\operatorname{Pic}^0_{V/k,\mathrm{red}}
                       \longrightarrow J_{Y_{n-1}}
\]
therefore has finite kernel. Formula (MA.5) gives
\begin{equation}
 (b_1\cdots b_{n-1})\,
       \mu_{L_1,\ldots,L_{n-1}}=r^\vee\lambda_{Y_{n-1}}r.
                                                        \tag{MA.7}
\end{equation}
The right side is a polarization by finite pullback of the
theta line bundle
\citep[Lemma 30.17.2 and its forward proof]{frC3beStacksFiniteAmple}.
When $n=1$, there are no cuts, the product is one, and (MA.7)
is the principal polarization of the curve.
In particular the argument is valid in characteristic $2$ and
does not rely on separability of the resulting isogenies.

We will also use $\mu_{M,H,\ldots,H}$ when $M$ is arbitrary
and $H$ is ample. This is defined by the one-cycle construction,
and is additive in $M$. It is symmetric: express $M$ as a
difference of two ample line-bundle classes and apply (MA.7)
to both. This avoids assigning a polarization to a nonample
divisor itself.

\emph{2. The exact mixed weighted identity.}
Let $C_e$ run through the connected components of $E_i\cap E_j$,
$i\ne j$, with an orientation $e=(i,j)$.
Each $C_e$ is smooth projective of dimension $d-1$.
Set
\[
 A_e=\operatorname{Pic}^0_{C_e/k,\mathrm{red}},\qquad
 \nu_e=\mu_{\,\underbrace{H|_{C_e},\ldots,H|_{C_e}}_{d-2}}
                            :A_e\longrightarrow A_e^\vee.
\]
For $d=2$, this means $\nu_e=\lambda_{C_e}$.
For $d>2$, (MA.7) says that a positive multiple of $\nu_e$
is a polarization. Let $r_{ie}:A_i\to A_e$ and
$r_{je}:A_j\to A_e$ be the two restrictions. Put
$N_i=\mathcal O_Z(E_i)|_{E_i}$.
The global choice of $H$ means that its restrictions from
both sides of every $C_e$ coincide. Projection and
self-intersection give the blocks of (MA.1):
\begin{align}
 (\Lambda_H)_{ii}
       &=-\mu_{\,N_i,H_i,\ldots,H_i},\notag\\
 (\Lambda_H)_{ij}
       &=-\sum_{e=(i,j)}r_{ie}^\vee\nu_e r_{je}
                                     \qquad(i\ne j).   \tag{MA.8}
\end{align}
This uses the one-cycle
$c_1(L_i)c_1(H_i)^{d-2}$ on each $E_i$.
In the off-diagonal term, its restriction to $C_e$ is the
zero-cycle $c_1(L_i|_{C_e})c_1(H|_{C_e})^{d-2}$.
The Albanese pushforward is dual to Picard restriction,
which gives exactly (MA.8).

Let $T=\operatorname{diag}([a_i])$ and
$d_e:\prod_iA_i\to A_e$ be
$d_e(x)=r_{ie}(x_i)-r_{je}(x_j)$.
The equality of line-bundle classes
\[
 D|_{E_i}=a_iN_i+
             \sum_{j\ne i}\sum_{e=(i,j)}a_jC_e
\]
and multilinearity of the one-cycle intersection imply
\begin{equation}
 T^\vee\Lambda_HT
 =\operatorname{diag}
    \bigl(a_i\mu_{\,\mathcal O_{E_i}(-D),H_i,\ldots,H_i}\bigr)
       +\sum_{e=(i,j)}a_i a_j\,d_e^\vee\nu_ed_e.
                                                        \tag{MA.9}
\end{equation}
This is the same block identity as AM.8; the new polarization
on a double stratum replaces its curve principal polarization.
Triple and higher strata do not add terms to this identity.

Choose the complete-intersection curves of Step~1 for the
positive diagonal terms and for each $\nu_e$, and take a common
positive multiple clearing all their integers in (MA.7).
Then that multiple of (MA.9) is the polarization of a product
theta line bundle pulled back under a map to a product of
Jacobians. Its diagonal components already have finite kernel.
The map is therefore finite onto its image, and its theta
pullback is ample. Thus $\Lambda_H$ is an isogeny, proving
the intrinsic assertion in every characteristic.

\emph{3. Remove the degree coordinates in the resolved setting.}
C.3bn gives the point comparison, the equality of ordinary and
perfect-family relations, and
\begin{equation}
 B^d(X)_{\mathbf Q}
 =S_d(E)/(J_K^d(E)+\theta_d\mathrm{CH}_1(E)_{\mathbf Q}).
                                                        \tag{MA.10}
\end{equation}
In an actual proper perfect family, the degree of its
codimension-$d$ Chern character on each component is constant.
Consequently $J_K^d(E)\subset S_{d,0}(E)$.

Let $c_j=c_1(H_j)^{d-1}\cap[E_j]\in\mathrm{CH}_1(E_j)$.
The degree vector of $\theta_d(c_j)$ is column $j$ of
\[
 M_{ij}=\deg\bigl(c_1(\mathcal O_Z(E_i)|_{E_j})
                              c_1(H_j)^{d-1}\bigr).
\]
This matrix is symmetric: for $i\ne j$ it is the degree of
$H^{d-1}$ on $E_i\cap E_j$. Its off-diagonal entries are
nonnegative. With $a=(a_i)$,
\[
 (Ma)_i=\deg\bigl(c_1(\mathcal O_{E_i}(D))
                                      c_1(H_i)^{d-1}\bigr)<0.
\]
The strict inequality follows by cutting with sufficiently
high powers of $H_i$ to a smooth curve and using the positive
degree of the ample $\mathcal O_{E_i}(-D)$ there.
For $A_0=\operatorname{diag}(a_i)$, the matrix $A_0MA_0$ has
nonnegative off-diagonal entries and strictly negative row sums.
The exact quadratic identity
\[
 x^{\mathsf T}(A_0MA_0)x
 =\sum_i\Bigl(\sum_j(A_0MA_0)_{ij}\Bigr)x_i^2
     -\sum_{i<j}(A_0MA_0)_{ij}(x_i-x_j)^2
\]
makes it negative definite. Hence $M$ is invertible over
$\mathbf Q$, and exceptional cycles span every component-degree
coordinate. Formula (MA.10) therefore becomes
\begin{equation}
 B^d(X)_{\mathbf Q}
 =S_{d,0}(E)/
 \bigl(J_K^d(E)+
       (\theta_d\mathrm{CH}_1(E)_{\mathbf Q}\cap S_{d,0}(E))\bigr).
                                                        \tag{MA.11}
\end{equation}

\emph{4. The Picard inputs are actual perfect-family relations.}
For $L_i\in A_i(k)$ choose a smooth connected projective curve
$P$ in $A_i$ through $0$ and $L_i$, omitting the trivial
zero-dimensional case. A normalized Poincare bundle
$\mathcal P_i$ restricts to an actual line bundle on
$E_i\times P$. The cycle
\[
 c_1(\mathcal P_i)\,
          \operatorname{pr}_{E_i}^*c_1(H_i)^{d-2}
                          \cap[E_i\times P]
\]
has codimension $d-1$ on that product. Proper pushforward
into the fixed $Z\times P$ gives an actual codimension-$d$
ordinary cycle family with endpoint difference
$j_{i*}z_i(L_i)$. The comparison AE.13 proves
\begin{equation}
                    \theta_d(z_i(L_i))\in J_K^d(E).     \tag{MA.12}
\end{equation}
If desired, this is obtained directly by taking the pure
weight-$d$ inverse Chern-character class of the displayed
cycle on a smooth finite model of $Z\times P$ and restricting
the resulting actual perfect class to $E\times P$.
Thus (MA.12) uses no perfect pushforward from a nonregular
component immersion into $E$ itself.

The component Albanese images of these classes are
$\Lambda_H(\prod_iA_i(k))$. Since an isogeny over an
algebraically closed field is surjective on points,
Step~2 proves (MA.4). These classes have degree zero on
each component because $L_i$ is algebraically trivial.

\emph{5. The exact denominator after both reductions.}
Set $Q_d=\theta_d\mathrm{CH}_1(E)_{\mathbf Q}$.
By (MA.4),
\[
 U_H\subset J_K^d(E)\cap Q_d\cap S_{d,0}(E),\qquad
 \alpha_E(U_H)=\bigoplus_iB_i(k)\otimes\mathbf Q.
\]
Every class in $S_{d,0}(E)$ can therefore be shifted by an
element of $U_H$ into $S_{d,00}(E)$.
If $j+q\in S_{d,00}(E)$, with $j\in J_K^d(E)$ and
$q\in Q_d\cap S_{d,0}(E)$, choose $u\in U_H$ with
$\alpha_E(u)=\alpha_E(q)$. Then
\[
 j+q=(j+u)+(q-u),\quad
 j+u\in J_K^d(E)\cap S_{d,00}(E),\quad
 q-u\in Q_d\cap S_{d,00}(E).
\]
The reverse containment is immediate. Applying this equality
to (MA.11) proves (MA.3).
\end{proof}

\emph{Exact scope and remaining question.}
The mixed polarization identity and its isogeny hold in every
characteristic for the stated smooth SNC components with their
negative ample combination. The identification (MA.3) with the
ordinary top quotient uses C.3bn and hence is asserted here only
in its positive-characteristic resolved algebraizable scope.
A non-SNC exceptional support satisfies C.3bn but need not admit
the component argument above. In dimensions where the required
resolution has not been supplied, this proposition does not
supply one. Even in the displayed scope, component Albanese-kernel
classes and the global compatibility of proper perfect families
across the SNC strata remain; neither is assumed to vanish.

